\documentclass[11pt,a4paper]{article}
\usepackage[utf8]{inputenc}
\usepackage[T1]{fontenc}
\usepackage{babel}
\renewcommand{\contentsname}{Table des mati\`eres}
\usepackage{lmodern}
\usepackage[margin=2cm]{geometry}
\usepackage{xcolor}
\usepackage{titlesec}
\usepackage{longtable}
\usepackage{booktabs}
\usepackage{colortbl}
\usepackage{graphicx}
\usepackage{parskip}
\usepackage{hyperref}
\usepackage{tikz}
\usepackage{xurl}
\urlstyle{tt}
\usetikzlibrary{shapes.geometric, arrows.meta, positioning}

\definecolor{navy}{HTML}{111111}
\definecolor{gold}{HTML}{1565C0}
\definecolor{greytext}{HTML}{555555}
\definecolor{lightgrey}{HTML}{F2F4F9}
\definecolor{okgreen}{HTML}{2E7D32}

\titleformat{\section}{\Large\bfseries\color{navy}}{\thesection}{0.6em}{}[{\color{navy}\titlerule}]
\titleformat{\subsection}{\large\bfseries\color{gold}}{\thesubsection}{0.6em}{}

\hypersetup{colorlinks=true, linkcolor=navy, urlcolor=navy, citecolor=navy}

\newcommand{\tblhead}[1]{\textcolor{white}{\textbf{#1}}}

\begin{document}

\begin{titlepage}
\centering
\vspace*{2cm}
{\Huge\bfseries\color{navy} Stockage, routage et choix techniques\par}
\vspace{0.5cm}
{\Large\bfseries\color{gold} R\'eponses aux questions de l'encadrante -- Syst\`eme RAG pour hcp.ma\par}
\vspace{1.5cm}
{\large\color{greytext} Structure d'accueil~: Haut-Commissariat au Plan (HCP), Direction des Syst\`emes d'Information Statistiques\par}
\vspace{0.3cm}
{\large\color{greytext} Bin\^ome~: Ayman El Baida et Saad Mesbah\par}
\vspace{0.3cm}
{\large\color{greytext} P\'eriode du stage~: 1er juillet -- fin ao\^ut 2026\par}
\vspace{0.3cm}
{\large\itshape\color{greytext} Note r\'edig\'ee suite \`a la r\'eunion du 17 juillet 2026 -- 20 juillet 2026\par}
\vfill
\end{titlepage}

\tableofcontents
\newpage

\section{Introduction}
Lors de la r\'eunion du 17 juillet 2026, quatre questions ont \'et\'e soulev\'ees sur l'architecture de stockage et de routage du syst\`eme~: (1) comment et o\`u sont stock\'ees les deux natures de donn\'ees produites par le pipeline d'ingestion (texte narratif et indicateurs chiffr\'es), et quelles en sont les limites~; (2) l'int\'er\^et d'introduire Redis~; (3) pourquoi diviser la base de donn\'ees en deux syst\`emes plut\^ot qu'un seul, le Routeur devant interroger chacun ind\'ependamment~; (4) comment le Routeur distingue une question chiffr\'ee d'une question narrative, et comment il g\`ere le cas o\`u une question rel\`eve des deux \`a la fois. Ce document r\'epond aux quatre points, puis pr\'esente le benchmarking demand\'e sur les choix techniques initiaux de l'ADR 0002 (\url{docs/adr/0002-stack-technique-prototype.md}).

\section{Stockage des donn\'ees apr\`es collecte}

\subsection{\'Etat actuel et limites}
Le pipeline produit deux natures de donn\'ees distinctes (ADR 0001, \url{docs/adr/0001-separation-texte-indicateurs.md}), stock\'ees dans deux syst\`emes diff\'erents~:

\begin{longtable}{@{}p{3.3cm}p{4.6cm}p{7.4cm}@{}}
\rowcolor{navy}
\tblhead{Donn\'ee} & \tblhead{Stockage actuel} & \tblhead{Limite connue} \\
\midrule
\endhead
Indicateurs chiffr\'es & Table \texttt{indicateur} (SQLite, \texttt{db/schema.sql}) & Fichier unique, un seul processus \'ecrivain \`a la fois (mode WAL permet des lectures concurrentes, pas des \'ecritures concurrentes). Suffisant pour un prototype \`a faible concurrence~; deviendrait un go\'ulot d'\'etranglement avec de nombreux utilisateurs simultan\'es en production. \\
Chunks de texte narratif & Pr\'evu~: ChromaDB (index vectoriel) + index BM25 en parall\`ele (ADR 0002) & ChromaDB local a un plafond pratique autour de quelques millions de vecteurs et filtre les m\'etadonn\'ees \emph{apr\`es} la recherche par similarit\'e (moins pr\'ecis \`a grande \'echelle qu'un filtrage pr\'ealable). Tr\`es adapt\'e au d\'eveloppement local et \`a un prototype, pas con\c{c}u pour plusieurs serveurs. \\
Fichiers bruts (PDF/XLSX/DOCX) & Disque, \texttt{data/raw/}, hors base de donn\'ees et hors d\'ep\^ot git (\texttt{.gitignore}) & Volume qui cro\^it sans limite ni sauvegarde automatique~; regenerable en relan\c{c}ant le Scraper, donc pas critique au stade actuel, mais \`a surveiller si le corpus grossit significativement. \\
\end{longtable}

Le principe qui gouverne ce choix~: chaque nature de donn\'ee va dans le syst\`eme optimis\'e pour son mode d'acc\`es. Les indicateurs se cherchent par \'egalit\'e exacte (\guillemotleft~donne-moi la ligne o\`u \texttt{code\_bds='I3181'} et \texttt{periode='2025T4'}\guillemotright) -- le cas d'usage historique d'une base relationnelle. Les chunks se cherchent par proximit\'e s\'emantique (\guillemotleft~trouve les passages qui parlent du m\^eme sujet que cette question\guillemotright) -- ce qu'une base relationnelle ne fait pas nativement, et ce pour quoi un index vectoriel existe.

\subsection{Redis comme couche de cache -- proposition}
L'ADR 0004 (\url{docs/adr/0004-api-bds-pour-les-indicateurs.md}, section \guillemotleft~pr\'ecisions du 15 juillet\guillemotright) d\'ecrit d\'ej\`a, sans le nommer, un sch\'ema \emph{cache-aside}~: pr\'e-remplissage nocturne des indicateurs cur\'es, appel direct \`a l'API BDS si absent du cache local, puis \'ecriture du r\'esultat pour enrichir ce cache. Redis est l'outil standard pour impl\'ementer exactement ce sch\'ema~: une base cl\'e-valeur en m\'emoire, donc des lectures nettement plus rapides qu'une base sur disque comme SQLite, avec une expiration automatique par cl\'e (TTL) -- pertinent ici puisque les indicateurs BDS se mettent \`a jour \`a un rythme connu (mensuel pour l'IPC et l'IPPI, par exemple), ce qui permet de fixer une dur\'ee de vie de cache align\'ee sur cette fr\'equence plut\^ot que de coder cette logique de p\'eremption \`a la main.

Redis ne remplacerait pas SQLite mais viendrait \emph{devant}~: SQLite reste la source de v\'erit\'e durable, celle qui porte la tra\c{c}abilit\'e (chaque indicateur reli\'e \`a un \texttt{id\_document}, exigence NF3 du cadrage) -- un r\^ole que Redis ne remplit pas nativement, \'etant volatile par d\'efaut. Redis jouerait le r\^ole d'une couche rapide pour les indicateurs consult\'es fr\'equemment, avec SQLite comme filet de s\'ecurit\'e derri\`ere pour toute donn\'ee absente du cache m\'emoire. C'est une \'evolution propos\'ee, pas encore impl\'ement\'ee~; elle fera l'objet d'un ADR d\'edi\'e (ADR 0006) si elle est retenue.

\section{Pourquoi diviser la base de donn\'ees}
Deux arguments distincts justifient cette s\'eparation, l'un fonctionnel, l'autre technique.

\textbf{Argument fonctionnel (ADR 0001).} Un institut statistique officiel ne peut pas se permettre qu'une recherche par similarit\'e renvoie un chiffre approximatif \`a la place du chiffre officiel exact. Les indicateurs doivent donc passer par une requ\^ete exacte, jamais par une recherche floue -- ce qui exclut d'office de les stocker uniquement dans un index vectoriel.

\textbf{Argument technique (ad\'equation outil/usage).} M\^eme sans la contrainte d'exactitude, ce sont deux probl\`emes de recherche diff\'erents, chacun avec son outil de pr\'edilection. Une base relationnelle excelle \`a r\'epondre \`a une requ\^ete exacte et index\'ee~; elle n'est pas con\c{c}ue pour comparer des vecteurs de similarit\'e. Un index vectoriel excelle \`a trouver des contenus s\'emantiquement proches~; il n'est pas con\c{c}u pour des requ\^etes structur\'ees exactes. Utiliser deux syst\`emes, chacun pour ce qu'il fait le mieux, est un principe connu sous le nom de \emph{polyglot persistence}. Le co\^ut~: deux syst\`emes \`a maintenir plut\^ot qu'un seul. Le b\'en\'efice~: chacun fait bien son travail plut\^ot qu'un seul qui ferait les deux moyennement bien.

\section{Le Routeur~: distinguer chiffr\'e et narratif, et g\'erer le cas mixte}

\subsection{Limite du mod\`ele actuel}
\texttt{src/routeur.py} est aujourd'hui un squelette~: une \'enum\'eration \texttt{TypeQuestion} \`a deux valeurs mutuellement exclusives, \texttt{CHIFFRE} ou \texttt{NOTION}, et une m\'ethode \texttt{classifier(question)} non encore impl\'ement\'ee (pr\'evue semaine 3, classification par un prompt LLM court). Ce mod\`ele \`a choix unique ne g\`ere pas correctement une question qui rel\`eve des deux \`a la fois -- exactement le cas soulev\'e en r\'eunion. Une question comme \guillemotleft~pourquoi le ch\^omage a-t-il augment\'e ces derniers trimestres~?\guillemotright~a une composante chiffr\'ee (le taux, par p\'eriode) et une composante narrative (les causes, pr\'esentes dans les rapports). Forcer un choix unique conduit \`a mal traiter ce genre de question, quel que soit le choix fait.

\subsection{Proposition r\'evis\'ee~: deux signaux ind\'ependants}
La r\'evision propos\'ee remplace le choix unique par deux signaux bool\'eens ind\'ependants -- \guillemotleft~cette question a-t-elle une composante chiffr\'ee~?\guillemotright\ et \guillemotleft~a-t-elle une composante narrative~?\guillemotright\ -- les deux pouvant \^etre vrais simultan\'ement.

\begin{longtable}{@{}p{4.2cm}p{11.9cm}@{}}
\rowcolor{navy}
\tblhead{\'Etape} & \tblhead{Fonctionnement propos\'e} \\
\midrule
\endhead
Signal chiffr\'e & Comparer la question aux 832 libell\'es du catalogue BDS d\'ej\`a disponible (\url{data/bds_catalogue.json}), par correspondance de mots-cl\'es ou de similarit\'e. Signal peu co\^uteux \`a calculer, r\'eutilise une donn\'ee d\'ej\`a en place. \\
Signal narratif & D\'etection d'une intention explicative -- mots d\'eclencheurs (\guillemotleft~pourquoi\guillemotright, \guillemotleft~comment\guillemotright, \guillemotleft~expliquer\guillemotright, \guillemotleft~tendance\guillemotright) ou, en V2, un prompt LLM court d\'edi\'e. \\
Dispatch & Si signal chiffr\'e seul~: \texttt{LookupStructure} uniquement. Si signal narratif seul~: \texttt{RetrievalReranker} uniquement. Si les deux~: les deux modules sont interrog\'es en parall\`ele. \\
Fusion & \texttt{Generateur} compose la r\'eponse \`a partir de ce qui a \'et\'e trouv\'e -- le chiffre exact d'un c\^ot\'e, l'explication sourc\'ee de l'autre -- avec citation des deux sources si les deux chemins ont produit un r\'esultat. \\
\end{longtable}

Cette r\'evision n'invalide pas le travail d\'ej\`a fait~: \texttt{LookupStructure} et \texttt{RetrievalReranker} restent chacun responsables d'un seul type de recherche (ADR 0001 inchang\'e). Ce qui change, c'est uniquement la couche de d\'ecision au-dessus -- le Routeur peut d\'esormais activer les deux chemins pour une m\^eme question plut\^ot que de choisir l'un ou l'autre. Cette r\'evision sera document\'ee formellement d\`es le d\'ebut de son impl\'ementation (semaine 3, Sprint 3).

\newpage
\section{Benchmarking des choix techniques (ADR 0002)}
Comparatif demand\'e des choix initiaux face aux alternatives courantes, sur la base de comparatifs ind\'ependants r\'ecents et, quand elle existe, de l'exp\'erience r\'eelle d\'ej\`a accumul\'ee sur ce projet.

\subsection{Scraping~: \texttt{requests} + \texttt{BeautifulSoup}}
\begin{longtable}{@{}p{3.6cm}p{11.9cm}@{}}
\rowcolor{navy}
\tblhead{Option} & \tblhead{\'Evaluation} \\
\midrule
\endhead
\texttt{requests} + \texttt{BeautifulSoup} (retenu) & L\'eger, rapide \`a mettre en \oe uvre, suffisant tant que les pages sont du HTML statique g\'en\'er\'e c\^ot\'e serveur -- confirm\'e empiriquement~: les 27 pages test\'ees le 15 juillet se sont r\'ecup\'er\'ees et pars\'ees sans ex\'ecution JavaScript n\'ecessaire. \\
Scrapy & Framework complet (concurrence, files de reprise, pipelines de traitement int\'egr\'es). Justifi\'e pour des milliers de pages avec crawl r\'ecurrent programm\'e~; surdimensionn\'e pour les quelques dizaines de pages cibl\'ees par ce projet. \\
Playwright / Selenium & N\'ecessaire uniquement si le contenu est g\'en\'er\'e c\^ot\'e client par JavaScript (site \guillemotleft~dynamique\guillemotright). Non justifi\'e ici~: v\'erifi\'e en conditions r\'eelles que hcp.ma sert du HTML directement exploitable. \\
\end{longtable}

\subsection{Extraction PDF~: \texttt{pdfplumber}}
\begin{longtable}{@{}p{3.6cm}p{11.9cm}@{}}
\rowcolor{navy}
\tblhead{Option} & \tblhead{\'Evaluation} \\
\midrule
\endhead
\texttt{pdfplumber} (retenu) & Contr\^ole fin de la d\'etection de tableaux, adapt\'e aux mises en page ambigu\"es -- pr\'ecis\'ement le cas des rapports hcp.ma. \textbf{D\'ej\`a valid\'e en conditions r\'eelles~: 30 PDF authentiques trait\'es sans erreur le 17 juillet 2026, de 2 \`a 645 tableaux extraits selon le fichier.} Cette preuve empirique p\`ese plus que n'importe quel comparatif th\'eorique. \\
PyMuPDF (\texttt{fitz}) & Le plus rapide pour l'extraction de texte brut \`a grande \'echelle~; d\'etection de tableaux plus r\'ecente et moins fine que \texttt{pdfplumber} sur des mises en page complexes. Utilis\'e en \emph{repli} dans l'ADR 0002 si \texttt{pdfplumber} montre ses limites de performance (les 2 fichiers volumineux \guillemotleft~Tabulations\guillemotright\ identifi\'es comme lents). \\
Camelot & Meilleure pr\'ecision de d\'etection sur des tableaux tr\`es structur\'es (factures, rapports \`a grille nette), mais d\'ependance suppl\'ementaire (Ghostscript) et moins robuste sur des mises en page vari\'ees comme celles observ\'ees sur hcp.ma. \\
\end{longtable}

\subsection{Embeddings~: BGE-M3}
\begin{longtable}{@{}p{3.6cm}p{11.9cm}@{}}
\rowcolor{navy}
\tblhead{Option} & \tblhead{\'Evaluation} \\
\midrule
\endhead
BGE-M3 (retenu) & Multilingue (plus de 100 langues, dont fran\c{c}ais et arabe -- pertinent si le support de l'arabe est un jour activ\'e), open source (licence MIT, h\'ebergeable localement, aucun co\^ut par appel ni d\'ependance \`a une API externe). Combine recherche dense, parcimonieuse et multi-vecteurs (ColBERT) dans un seul mod\`ele, avec un contexte de 8~000 tokens. Sur les classements MTEB r\'ecents, l\'eg\`erement en retrait par rapport \`a OpenAI text-embedding-3-large sur les t\^aches uniquement anglaises, mais nettement au-dessus sur les t\^aches multilingues. \\
OpenAI text-embedding-3-large & Tr\`es bon score MTEB, infrastructure g\'er\'ee sans installation. Mais~: API payante par appel, d\'ependance externe (pertinent au regard du point encore ouvert de l'ADR 0002 sur la politique de s\'ecurit\'e des donn\'ees du HCP), et les documents restent envoy\'es \`a un tiers externe. \\
multilingual-e5-large & Alternative open source multilingue s\'erieuse, performances proches de BGE-M3 sur plusieurs t\^aches, mais sans le support natif dense+parcimonieux+multi-vecteurs qui simplifie l'impl\'ementation de la recherche hybride pr\'evue. \\
\end{longtable}

\subsection{Index vectoriel~: ChromaDB}
\begin{longtable}{@{}p{3.6cm}p{11.9cm}@{}}
\rowcolor{navy}
\tblhead{Option} & \tblhead{\'Evaluation} \\
\midrule
\endhead
ChromaDB (retenu) & Local, installation quasi instantan\'ee, id\'eal pour le d\'eveloppement et un prototype -- plafond pratique autour de quelques millions de vecteurs, tr\`es au-dessus du volume attendu pour ce projet (corpus de quelques dizaines \`a quelques centaines de documents). \\
FAISS & Le plus rapide en recherche brute, mais c'est une biblioth\`eque de calcul, pas une base de donn\'ees~: pas de persistance ni de gestion des m\'etadonn\'ees int\'egr\'ee, \`a construire soi-m\^eme par-dessus. \\
Qdrant & Con\c{c}u pour la production~: filtrage des m\'etadonn\'ees appliqu\'e \emph{avant} la recherche (plus correct \`a grande \'echelle), scalable horizontalement. Pertinent si HCP d\'eploie ce syst\`eme en production avec un corpus en forte croissance -- chemin de migration naturel depuis ChromaDB si ce cap est atteint. \\
pgvector & Extension PostgreSQL~: pas de nouveau syst\`eme si une base PostgreSQL est d\'ej\`a en place. \`A envisager si l'infrastructure HCP standardise sur PostgreSQL plut\^ot que sur des bases sp\'ecialis\'ees. \\
\end{longtable}

\subsection{Recherche lexicale~: \texttt{rank-bm25}}
\begin{longtable}{@{}p{3.6cm}p{11.9cm}@{}}
\rowcolor{navy}
\tblhead{Option} & \tblhead{\'Evaluation} \\
\midrule
\endhead
\texttt{rank-bm25} (retenu) & Biblioth\`eque Python pure (d\'ependance~: NumPy uniquement), aucun serveur \`a installer ou maintenir. Adapt\'e aux corpus de moins de 100~000 documents -- tr\`es au-dessus du volume de ce projet. \\
BM25S & \'Evolution r\'ecente de \texttt{rank-bm25}, m\^eme simplicit\'e d'usage mais nettement plus rapide (calcul matriciel creux). \`A consid\'erer comme mise \`a niveau directe si la vitesse de recherche devenait un jour un point sensible -- aucun changement d'architecture requis. \\
Elasticsearch / OpenSearch & Standard en production \`a grande \'echelle (millions de documents, plusieurs n\oe uds), mais n\'ecessite un serveur Java d\'edi\'e et une charge op\'erationnelle largement hors de proportion avec le corpus actuel. \\
\end{longtable}

\subsection{Base des indicateurs~: SQLite}
\begin{longtable}{@{}p{3.6cm}p{11.9cm}@{}}
\rowcolor{navy}
\tblhead{Option} & \tblhead{\'Evaluation} \\
\midrule
\endhead
SQLite (retenu) & Conforme ACID, mode WAL permettant des lectures concurrentes, z\'ero maintenance (pas de serveur \`a d\'eployer), tr\`es rapide pour des recherches ponctuelles par cl\'e exacte -- exactement le sc\'enario de \texttt{LookupStructure}. Limite~: un seul \'ecrivain \`a la fois, adapt\'e \`a une faible concurrence. \\
PostgreSQL & Supporte de nombreux \'ecrivains simultan\'es, choix par d\'efaut pour une application multi-utilisateurs en production. \`A envisager pour un d\'eploiement r\'eel \`a HCP au-del\`a du prototype. \\
DuckDB & Optimis\'e pour l'analytique sur de gros volumes (requ\^etes d'agr\'egation, colonnes plut\^ot que lignes) -- ne correspond pas au sc\'enario ici, qui est de la recherche ponctuelle de lignes individuelles, pas de l'analyse en masse. \\
\end{longtable}

\subsection{Interface de d\'emonstration~: Streamlit}
\begin{longtable}{@{}p{3.6cm}p{11.9cm}@{}}
\rowcolor{navy}
\tblhead{Option} & \tblhead{\'Evaluation} \\
\midrule
\endhead
Streamlit (retenu) & Plus flexible pour une interface multi-\'ecrans (question/r\'eponse, mais aussi visualisation d'indicateurs \'eventuelle) et une personnalisation pouss\'ee de la mise en page. \\
Gradio & Chemin le plus rapide pour un simple chatbot -- composants pr\^ets \`a l'emploi pour l'interaction conversationnelle. Pertinent si le besoin se limite strictement \`a un fil de discussion, sans tableau de bord annexe. \\
\end{longtable}

\section{Synth\`ese et prochaines \'etapes}
Les quatre points soulev\'es en r\'eunion aboutissent \`a trois actions concr\`etes~: (1) documenter formellement l'introduction \'eventuelle de Redis comme couche de cache devant SQLite, sous forme d'un ADR 0006~; (2) r\'eviser \texttt{routeur.py} pour remplacer le choix unique \texttt{CHIFFRE}/\texttt{NOTION} par deux signaux ind\'ependants capables de se combiner, avant le d\'ebut de son impl\'ementation pr\'evue en semaine 3~; (3) le choix des briques techniques actuelles (ADR 0002) reste justifi\'e au stade prototype pour chacune d'elles, avec pour chacune un chemin de montée en charge identifi\'e si HCP d\'ecidait d'un d\'eploiement en production au-del\`a de ce stage.

\end{document}
